%Daten zur Institution

%\input{Dozentdaten}

%\renewcommand{\fachbereich}{Fachbereich}

%\renewcommand{\dozent}{Prof. Dr. . }

%Klausurdaten

\renewcommand{\klausurgebiet}{ }

\renewcommand{\klausurtyp}{ }

\renewcommand{\klausurdatum}{. 20}

\klausurvorspann {\fachbereich} {\klausurdatum} {\dozent} {\klausurgebiet} {\klausurtyp}

%Daten für folgende Punktetabelle

\renewcommand{\aeins}{  3  }

\renewcommand{\azwei}{  3   }

\renewcommand{\adrei}{  6  }

\renewcommand{\avier}{  4  }

\renewcommand{\afuenf}{  5  }

\renewcommand{\asechs}{  0  }

\renewcommand{\asieben}{  5  }

\renewcommand{\aacht}{  6  }

\renewcommand{\aneun}{  6  }

\renewcommand{\azehn}{  0  }

\renewcommand{\aelf}{  6  }

\renewcommand{\azwoelf}{  12  }

\renewcommand{\adreizehn}{  0   }

\renewcommand{\avierzehn}{  9  }

\renewcommand{\afuenfzehn}{  65  }

\renewcommand{\asechzehn}{    }

\renewcommand{\asiebzehn}{    }

\renewcommand{\aachtzehn}{    }

\renewcommand{\aneunzehn}{    }

\renewcommand{\azwanzig}{    }

\renewcommand{\aeinundzwanzig}{    }

\renewcommand{\azweiundzwanzig}{    }

\renewcommand{\adreiundzwanzig}{    }

\renewcommand{\avierundzwanzig}{    }

\renewcommand{\afuenfundzwanzig}{    }

\renewcommand{\asechsundzwanzig}{    }

\punktetabellevierzehn

\klausurnote

\newpage

\setcounter{section}{0}

\inputaufgabegibtloesung
{3}
{

Definiere die folgenden
\zusatzklammer {kursiv gedruckten} {} {} Begriffe.
\aufzaehlungsechs{Eine
\stichwort {Hauptkrümmung} {}
zu einer differenzierbaren Fläche

\mavergleichskette
{\vergleichskette
{ Y
}
{ \subseteq }{ \R^3
}
{  }{
}
{    }{
}
{   }{
}
}
{}{}{}
in einem Punkt

\mavergleichskette
{\vergleichskette
{ P
}
{ \in }{ Y
}
{  }{
}
{    }{
}
{   }{
}
}
{}{}{.}

}{Ein
\stichwort {Diffeomorphismus} {}
zwischen differenzierbaren Mannigfaltigkeiten
\mathkor {} {L} {und} {M} {.}

}{Ein
\stichwort {Tangentialvektor} {}
in einem Punkt

\mavergleichskette
{\vergleichskette
{  P
}
{ \in }{  M
}
{  }{
}
{    }{
}
{   }{
}
}
{}{}{}
einer differenzierbaren Mannigfaltigkeit $M$.

}{Ein
\stichwort {Differentialoperator erster Ordnung} {}
auf einer differenzierbaren Mannigfaltigkeit $M$.

}{Ein
\stichwort {euklidischer Halbraum} {.}

}{Eine
\definitionsverweis {kompakte Ausschöpfung}{}{}
eines
\definitionsverweis {topologischen Raumes}{}{}
$X$.
}

}
{} {}

\inputaufgabegibtloesung
{3}
{

Formuliere die folgenden Sätze.
\aufzaehlungdrei{Der Satz über den
\definitionsverweis {Tangentialraum}{}{}
einer
\definitionsverweis {abgeschlossenen Untermannigfaltigkeit}{}{}

\mavergleichskette
{\vergleichskette
{  M
}
{ \subseteq }{  N
}
{  }{
}
{    }{
}
{   }{
}
}
{}{}{.}}{Der Satz über die Gestalt einer zurückgezogenen Differentialform $\varphi^* \omega$  unter einer differenzierbaren Abbildung
\maabbdisp {\varphi} { U } { V
} {}
\zusatzklammer {mit
\definitionsverweis {offenen Teilmengen}{}{}
\mathkor {} {U \subseteq \R^n} {und} {V \subseteq \R^m} {,}
deren Koordinaten mit
\mathl{x_1  , \ldots , x_n}{} bzw. mit
\mathl{y_1  , \ldots , y_m}{} bezeichnet seien} {} {}
von einer
$k$-\definitionsverweis {Differentialform}{}{}
auf $V$ mit der Darstellung

\mavergleichskettedisp
{\vergleichskette
{  \omega
}
{ =} {\sum_{I,\,  { \# \left(   I  \right) } = k}  f_I  dy_I
}
{ } {
}
{ } {
}
{ } {
}
}
{}{}{,}
wobei
\maabb {f_I} { V } {\R
} {}
\definitionsverweis {Funktionen}{}{}
sind.}{Der Satz über die Charakterisierung von geodätischen Kurven bezüglich des Levi-Civita-Zusammenhangs auf einer riemannschen Mannigfaltigkeit $M$.}

}
{} {}

\inputaufgabegibtloesung
{6 (2+4)}
{

Es seien

\mavergleichskette
{\vergleichskette
{ \alpha,\beta
}
{ \in }{ \R_+
}
{  }{
}
{    }{
}
{   }{
}
}
{}{}{}
und sei

\mavergleichskettedisp
{\vergleichskette
{ Y
}
{ =} { \left\{ (x,y)\in\R^2  \mid   \alpha x^2+\beta y^2 = 1        \right\}
}
{ } {
}
{ } {
}
{ } {
}
}
{}{}{.}
\aufzaehlungzwei {Bestimme die
\definitionsverweis {Gauß-Abbildung}{}{}
zu $Y$.
} { Man gebe zur Gauß-Abbildung zu $Y$ explizit eine
\definitionsverweis {Umkehrabbildung}{}{}
an.
}

}
{} {}

\inputaufgabe
{4}
{

Man erläutere das Begriffspaar \anfuehrung{extrinsisch und intrinsisch}{} im Kontext von Mannigfaltigkeiten.

}
{} {}

\inputaufgabegibtloesung
{5}
{

Beweise die Formel für den Flächeninhalt der Oberfläche eines Rotationskörpers zu einer
\definitionsverweis {differenzierbaren Kurve}{}{}
\maabbeledisp {\gamma} {]a,b[} {\R^2
} { t } {(x(t),y(t))
} {,}
mit

\mavergleichskette
{\vergleichskette
{ y(t)
}
{ > }{ 0
}
{  }{
}
{    }{
}
{   }{
}
}
{}{}{.}

}
{} {}

\inputaufgabe
{0}
{

}
{} {}

\inputaufgabegibtloesung
{5}
{

Es sei $X$ ein
\definitionsverweis {Torus}{}{.}
Man gebe eine surjektive
\definitionsverweis {differenzierbare Abbildung}{}{}
\maabbdisp {\varphi} {\R^2} {X
} {}
derart an, dass auch die
\definitionsverweis {Tangentialabbildung}{}{}
\maabbdisp {T_P(\varphi)} { T_P\R^2 } { T_{\varphi(P)}X
} {}
in jedem Punkt

\mavergleichskette
{\vergleichskette
{ P
}
{ \in }{ \R^2
}
{  }{
}
{    }{
}
{   }{
}
}
{}{}{}
surjektiv ist.

}
{} {}

\inputaufgabegibtloesung
{6 (2+2+2)}
{

Wir betrachten den Graph $M$ der Abbildung
\maabbeledisp {\varphi} {\R^2} {\R
} {(u,v)} { u^2+uv-v^3
} {,} als zweidimensionale abgeschlossene Untermannigfaltigkeit des $\R^3$, also

\mavergleichskettedisp
{\vergleichskette
{ M
}
{ =} { \left\{ (u,v,u^2+uv-v^3)  \mid  (u,v) \in \R^2         \right\}
}
{ } {
}
{ } {
}
{ } {
}
}
{}{}{}
mit der vom $\R^3$ induzierten riemannschen Metrik. Es sei
\maabbeledisp {\psi} {\R^2} { M
} {(u,v)} { (u,v,u^2+uv-v^3)
} {,}
die zugehörige Diffeomorphie.

a) Bestimme das totale Differential zu $\psi$ sowie die Bildvektoren
\mathl{T_P(\psi) (e_1)}{} und
\mathl{T_P(\psi) (e_2)}{} in
\mathl{T_{\psi(P)}M}{.}

b) Bestimme für jeden Punkt der Form

\mavergleichskette
{\vergleichskette
{ P
}
{ = }{ (u,0)
}
{  }{
}
{    }{
}
{   }{
}
}
{}{}{}
den Flächeninhalt des von
\mathl{T_P(\psi) (e_1)}{} und
\mathl{T_P(\psi) (e_2)}{} in
\mathl{T_{\psi(P)}M}{} aufgespannten Parallelogramms.

c) Bestimme für jeden Punkt der Form

\mavergleichskette
{\vergleichskette
{ P
}
{ = }{ (0,v)
}
{  }{
}
{    }{
}
{   }{
}
}
{}{}{}
den Flächeninhalt des von
\mathl{T_P(\psi) (e_1)}{} und
\mathl{T_P(\psi) (e_2)}{} in
\mathl{T_{\psi(P)}M}{} aufgespannten Parallelogramms.

}
{} {}

\inputaufgabegibtloesung
{6}
{

Beweise die Produktregel für differenzierbare Differentialformen auf einer offenen Menge des $\R^n$.

}
{} {}

\inputaufgabe
{0}
{

}
{} {}

\inputaufgabegibtloesung
{6}
{

Es sei $M$ eine
\definitionsverweis {Mannigfaltigkeit mit Rand}{}{}
und sei

\mavergleichskette
{\vergleichskette
{ P
}
{ \in }{ \partial M
}
{  }{
}
{    }{
}
{   }{
}
}
{}{}{}
ein Randpunkt. Zeige, dass für einen Tangentialvektor

\mavergleichskette
{\vergleichskette
{ v
}
{ \in }{ T_PM
}
{  }{
}
{    }{
}
{   }{
}
}
{}{}{}
folgende Eigenschaften äquivalent sind.
\aufzaehlungzwei {Es gibt einen stetig differenzierbaren Weg
\maabb {\gamma} { [- \epsilon, 0]} { M
} {}
mit

\mavergleichskette
{\vergleichskette
{ \gamma(0)
}
{ = }{ P
}
{  }{
}
{    }{
}
{   }{
}
}
{}{}{}
und

\mavergleichskette
{\vergleichskette
{ \gamma'(0)
}
{ = }{ v
}
{  }{
}
{    }{
}
{   }{
}
}
{}{}{.}
} { $v$ wird bei jeder Karte mit einem negativen Halbraum unter der Tangentialabbildung auf den rechten Halbraum abgebildet.
}

}
{} {}

\inputaufgabegibtloesung
{12}
{

Beweise den Satz von Stokes für Mannigfaltigkeiten mit Rand.

}
{} {}

\inputaufgabe
{0}
{

}
{} {}

\inputaufgabegibtloesung
{9 (1+1+3+2+2)}
{

Wir betrachten die trigonometrische Parametrisierung
\maabbeledisp {\psi} {\R } { S^1
} { t } { \left(  \cos  t   , \,   \sin  t                   \right)
} {}
des Einheitskreises. Es sei

\mavergleichskette
{\vergleichskette
{ c
}
{ \in }{ [0, 2 \pi[
}
{  }{
}
{    }{
}
{   }{
}
}
{}{}{}
fixiert. Auf dem trivialen Vektorbündel
\maabbdisp {} {S^1 \times \R^2} { S^1
} {}
sei ein
\definitionsverweis {linearer Zusammenhang}{}{}
$\nabla$ gegeben derart, dass der längs $\psi$
\definitionsverweis {zurückgezogene Zusammenhang}{}{}
$\psi^* \nabla$ durch die
\definitionsverweis {Christoffelsymbole}{}{}
\zusatzklammer {der Index für die einzige Ableitungsrichtung wird weggelassen} {} {}

\mavergleichskettedisp
{\vergleichskette
{ \begin{pmatrix} \tilde{\Gamma}_1^1 (t)   &  \tilde{\Gamma}_1^2 (t)   \\ \tilde{\Gamma}_2^1(t)  &  \tilde{\Gamma}_2^2(t)  \end{pmatrix}
}
{ =} { \begin{pmatrix}  0  &   - { \frac{  c  }{  2 \pi   } }   \\  { \frac{  c  }{  2 \pi  } }  &  0  \end{pmatrix}
}
{ } {
}
{ } {
}
{ } {
}
}
{}{}{}
gegeben ist. Wir betrachten zu einem Punkt

\mavergleichskette
{\vergleichskette
{ Q
}
{ = }{ \begin{pmatrix}  a  \\b   \end{pmatrix}
}
{ \in }{ \R^2
}
{    }{
}
{   }{
}
}
{}{}{}
die Abbildung
\maabbeledisp {\Psi_{Q}} {\R} { S^1 \times \R^2
} { t } { \left(  \cos  t   , \,   \sin  t  , \,   M(t)  \begin{pmatrix}  a  \\b   \end{pmatrix}                 \right)
} {}
mit

\mavergleichskettedisp
{\vergleichskette
{ M(t)
}
{ =} { \begin{pmatrix}  \cos  { \frac{  ct  }{  2 \pi   } }   &   -  \sin  { \frac{  ct  }{  2 \pi   } }    \\  \sin  { \frac{  ct  }{  2 \pi  } }   &  \cos  { \frac{  ct  }{  2 \pi   } }  \end{pmatrix}
}
{ } {
}
{ } {
}
{ } {
}
}
{}{}{.}
\aufzaehlungfuenf{Bestimme $\Psi_Q(0)$.
}{Bestimme $\Psi_Q(2 \pi)$.
}{Zeige, dass  $\Psi_Q$ eine
\definitionsverweis {horizontale Liftung}{}{}
längs $\psi$  ist.
}{Zeige, dass
\maabbeledisp {} {\R^2} { \R^2
} { Q } { \Psi_Q(2 \pi)
} {,}
eine
\definitionsverweis {lineare Isometrie}{}{}
ist
\zusatzklammer {der Basispunkt $(1,0)$ wird hier nicht aufgeführt} {} {.}
}{Zeige, dass
\maabbeledisp {} {\Z} { \operatorname{SL}_{  2  } \! \left(  \R  \right)
} { k } { \left(  Q \mapsto \Psi_Q(2k \pi )  \right)
} {,}
ein
\definitionsverweis {Gruppenhomomorphismus}{}{}
ist.
}

}
{} {}

 [[Kategorie:Latexseite]]
